\subsection{Overview} \label{sec:overview}
The primary purpose of \toolname is to accelerate pairwise sequence alignment through (1)~a memory-frugal and efficient algorithm, and (2)~optimized CPU, GPU, and ASIC implementations.

\toolname solves the \defn{approximate string matching (ASM)} problem with the \defn{edit distance}~\citep{levenshtein1966binary} as the cost metric. That is, given two strings, \var{text} and \var{pattern}, \toolname finds the minimum number of single-letter substitutions, insertions, and deletions to convert \var{text} into \var{pattern}. Additionally, the sequence of edits that corresponds the edit distance is reported, which is called \defn{CIGAR string}.

The \toolname algorithm is based on the GenASM algorithm~(\secref{genasm_algo}).
\citep{cali2020genasm} first proposed the GenASM algorithm as an algorithm/hardware co-design targeted for an ASIC accelerator, and demonstrated GenASM's potential for very high throughput and resource efficiency. However, as we show in \secref{challenges}, the GenASM algorithm (1)~requires large amounts of memory bandwidth, (2)~exhibits a large memory footprint, and (3)~does some unnecessary work. These inefficiencies limit GenASM's throughput and resource efficiency on both commodity and custom hardware, and addressing them is critical.

To this end, we propose \toolname's three novel algorithmic improvements to GenASM in \secref{optimizations}. 
In \secref{throughput}, we experimentally demonstrate that these improvements significantly increase performance on recent CPUs and GPUs. In \secref{sensitivity}, we explore the throughput behavior of GenASM with and without the proposed improvements across various configurations. We show in \secref{hw_benefits} that an ASIC implementation of \toolname will have significantly reduced chip area and power consumption compared to the ASIC designed for GenASM~\citep{cali2020genasm} while maintaining the same throughput. In \secref{accuracy}, we explore the accuracy behavior of GenASM and \toolname across various configurations.

\subsection{GenASM Algorithm} \label{sec:genasm_algo}
The GenASM algorithm~\citep{cali2020genasm} consists of two sub-algorithms: \emph{GenASM-DC} and \emph{GenASM-TB}. GenASM-DC~(\secref{bitap}) fills a bitvector-based dynamic programming table. The last column of the table indicates the edit distance between the two input strings. GenASM-TB (\secref{traceback}) re-traces this optimal solution in the constructed table. To better scale with longer input sequences, GenASM uses a \emph{windowing heuristic} (\secref{windowing}).

\subsubsection{GenASM-DC Algorithm} \label{sec:bitap}
GenASM-DC uses only cheap bitwise operations to calculate the edit distance between two strings \var{text} and \var{pattern}~\citep{cali2020genasm}.
It builds an (\var{n}+1)$\times$(\var{k}+1) dynamic programming (DP) table \var{R}, where \var{n}=$length$(\var{text}) and \var{k} is the maximum number of edits considered. The entries of \var{R} are \var{m}-bit bitvectors, where \var{m}=$length$(\var{pattern}). \figref{bitap_example} shows an example of \var{R} after it is constructed by GenASM-DC.

\begin{figure}[h]
    \centering
    \includegraphics[trim={8mm 0 0 0},width=0.87\columnwidth,keepaspectratio]{figures/bitap_example.pdf}
    \caption{An example of DP table \var{R} with \var{text}=\texttt{ACGT} and \var{k}=4. The bitmasks for \var{pattern}=\texttt{ACGA} are shown on the right. The colored arrows show the possible origins and data dependencies of the $0$ at \var{d}=1,\var{i}=2,\var{j}=2. The values in the red marked diagonal are mutually independent and thus can be computed in parallel.}
    \label{fig:bitap_example}
\end{figure}

\begin{theorem}
    \label{theorem:R_interpretation}
    The entries (bitvectors) of \var{R} can be interpreted as follows:
    \vspace{-5px}\begin{align*}
    \var{j}\text{-th bit of }\var{R}[\var{i}][\var{d}] = 0 \Longleftrightarrow \\[-5pt]
    distance\big(\var{text}[\var{i}:\var{n}), \var{pattern}[\var{j}:\var{m})\big) \leq \var{d}
    \end{align*}
\end{theorem}

In natural language, \theoref{R_interpretation} states that the \var{j}-th bit of the bitvector \var{R[i][d]} is 0 exactly if the suffix of \var{text} starting at character \var{i} and the suffix of \var{pattern} starting at character \var{j} differ by at most \var{d} edits. Following this interpretation, the first row \var{d=d$_{OPT}$} that has a \var{0} in the first bit (\var{j=0}) of the leftmost column (\var{i=0}) indicates that the edit distance between \var{text} and \var{pattern} is \var{d$_{OPT}$}.  This bit is marked in pink in~\figref{bitap_example}.

GenASM-DC~(\algref{bitap}) starts by pre-processing \var{pattern} into four \defn{pattern masks}, one per character in the alphabet. The pattern mask for character \var{X}{$\in$}$\{A,C,G,T\}$ is a bitvector of length \var{m}=$length$(\var{pattern}), with a 0 in the \var{i}-th bit if \var{pattern}[\var{i}]==\var{X}. See \figref{bitap_example} for an example.

GenASM-DC populates the rightmost column (\alglineref{right_init}) and topmost row (\alglineref{top_init}) of \var{R}. The remaining entries are then calculated from their respective neighbors in the north (\alglineref{insertion}, insertion), north-east (\algrangeref{deletion}{substitution}, deletion and substitution), and east (\alglineref{match}, match) through simple bitwise update rules. We refer to ~\citep{cali2020genasm,wu_manber_bitap,baeza-yates-gonnet-1992} for detailed arguments on the correctness of GenASM-DC.
To follow the rest of this paper, it is sufficient to consider (1)~the interpretation of \var{R} given in \theoref{R_interpretation}, and (2)~the north-east data dependencies imposed by \algref{bitap} and shown in \figref{bitap_example}.

\begin{algorithm}
\begin{spacing}{0.78}
\caption{GenASM-DC Algorithm}\label{alg:bitap}
\textbf{Inputs:} \var{text}, \var{pattern}, \var{k} \\
\textbf{Outputs:} \var{editDist}

\begin{algorithmic}[1]
    \State $\var{n} \gets \Call{length}{\var{text}}$
    \State $\var{m} \gets \Call{length}{\var{pattern}}$
    \State $\var{PM} \gets \Call{buildPatternMasks}{\var{pattern}}$
    \State
    \State $\var{R}[\var{n}][\var{d}] \gets 11...1 \ShiftLeft \var{d}$ \label{algline:right_init}
    \Comment{Initialize for all $0 \leq \var{d} \leq \var{k}$ }
    \State
    \For{\var{i} \textbf{in} $(\var{n}-1):-1:0$}
        \State $\var{char} \gets \var{text}[\var{i}]$
        \State $\var{curPM} \gets \var{PM}[\var{char}]$
        \State
        \State $\var{R}[\var{i}][0] \gets (\var{R}[\var{i}+1][0] \ShiftLeft 1) \BitOr \var{curPM}$ \label{algline:top_init}
        \Comment{exact match}
        \For{\var{d} \textbf{in} $1:\var{k}$}
            \State $\var{I} \gets \var{R}[\var{i}][\var{d}-1] \ShiftLeft 1$ \label{algline:update_rules_start}\label{algline:insertion}
            \Comment{insertion}
            \State $\var{D} \gets \var{R}[\var{i}+1][\var{d}-1]$\label{algline:deletion}
            \Comment{deletion}
            \State{$\var{S} \gets \var{R}[\var{i}+1][\var{d}-1] \ShiftLeft 1$}\label{algline:substitution}
            \Comment{substitution}
            \State{$\var{M} \gets (\var{R}[\var{i}+1][\var{d}] \ShiftLeft 1) \BitOr \var{curPM}$}\label{algline:match}
            \Comment{match}
            \State $\var{R}[\var{i}][\var{d}] \gets \var{I} \BitAnd \var{D} \BitAnd \var{S} \BitAnd \var{M}$ \label{algline:update_rules_end}
        \EndFor
    \EndFor
    \State
    \State $\var{editDist} \gets \argmin_\var{d}\{\Call{msb}{\var{R}[0][\var{d}]} = 0\}$ \label{algline:edit_dist}
\end{algorithmic}
\end{spacing}
\end{algorithm}

\noindent\textbf{Intra-Task Parallelism.} \label{sec:task_parallelism}
\citep{cali2020genasm} enables efficient intra-task parallelism by identifying that the DP entries within each north-west to south-east diagonal (one such diagonal is marked in red in \figref{bitap_example}) do \emph{not} depend on each other, hence they can be computed in parallel.

\subsubsection{GenASM-TB Algorithm} \label{sec:traceback}
For use-cases like read mapping, the pairwise sequence alignment algorithm should report both the edit distance and the corresponding sequence of edits, which is called the \defn{CIGAR string}. Obtaining the CIGAR string involves retracing the origin of the edit distance value as a linear path through DP entries in their reverse construction order; this process is called \defn{traceback}.


GenASM enables efficient traceback operations based on two key observations:
First, if \emph{all} intermediate values of variables \var{I}, \var{D}, \var{S} and \var{M} in \algref{bitap} are stored, then one can follow the path of \var{0}s in these variables, starting from \var{0} in the west of \var{R} that indicates the edit distance (highlighted in pink in \figref{bitap_example}) and go towards the north-east corner of \var{R}. Whenever a 0 in one of these variables is traversed, the name of that variable is recorded as an edit (e.g., '\var{I}' for an insertion). Second, it is sufficient to store only three out of the four variables (because \var{S} can be obtained by shifting \var{D}), saving both memory footprint and bandwidth. %

\subsubsection{GenASM's Windowing Heuristic} \label{sec:windowing}


To provide a linear runtime complexity, \citep{cali2020genasm} proposes a greedy \defn{windowing heuristic}. Instead of aligning \var{text} and \var{pattern} in a single run of GenASM-DC, the windowing heuristic runs GenASM-DC multiple times as a subroutine in \defn{windows} of size \var{W}. In each window, a prefix of size \var{W} characters of each sequence (i.e., $\var{text}[0:\var{W})$ and $\var{pattern}[0:\var{W})$) are aligned. The first $\var{W} - \var{O}$ characters of the window are greedily considered aligned optimally, where we call \var{O} the window \defn{overlap}. The smaller strings $\var{text}[\var{W}-\var{O}:\var{n})$ and $\var{pattern}[\var{W}-\var{O}:\var{m})$ then remain to be aligned in the next window.

This approach has three advantages. First, instead of constructing a large table of \var{n}$\times$\var{m}$\times$\var{k} bits, only $\frac{\var{m}}{\var{W} - \var{O}}$ tables of $\var{W}^3$ bits must be constructed, saving memory footprint, data movement, and computation. Second, the GenASM-DC subroutine now runs over constant-sized sequences, simplifying its implementation. For example, DP entries can be statically assigned to processing elements~\citep{cali2020genasm}, and the data movement and exact memory footprint are known at compile time, even if the length of the input sequences is unknown. Third, the program flow (e.g., the number of loop iterations per window) is entirely known at compile time, giving the compiler the ability to optimize.

The windowing strategy is greedy and heuristic, so it is possible that it could miss the optimal alignment and produce a sub-optimal one instead. This is a key limitation of GenASM and \toolname. Note that several state-of-the-art tools do not give any optimality guarantees either, and instead experimentally demonstrate their practical accuracy, as \toolname does. This includes greedy alignment techniques like SeGraM~\citep{cali2022segram}, Darwin~\citep{turakhia2019darwin}, and WFA-adaptive~\citep{wfa_algorithm}, as well as mappers based on sparse dynamic programming, like minimap2~\citep{ksw2_b}. To balance performance and accuracy, the tunable parameters \var{W}~(window size) and \var{O}~(window overlap) must be selected appropriately.
The parameter \var{W} can be understood as the \emph{range} of solutions considered, similar to the \defn{band width}~\citep{UKKONEN1985banded} in popular alignment implementations~(e.g. \citep{edlib,ksw2_a,ksw2_b}). The parameter \var{O} can be understood as the \emph{globality} of the solutions or \emph{inverse greediness}. We demonstrate in \secref{results} (1)~that higher \var{W} and \var{O} generally improve accuracy, at the cost of lowering throughput, (2)~that the best choice of \var{W} and \var{O} depends on the input dataset (e.g., its error distribution and read lengths), and (3)~that \var{W}=64 and \var{O}=33 achieve a good throughput/accuracy tradeoff for long and short read mapping.






\subsection{Inefficiencies in the GenASM Algorithm} \label{sec:challenges}
We identify three inefficiencies in the GenASM algorithm: (1)~it has a large amount of data movement, (2)~it has a large memory
footprint, and (3)~it does some unnecessary work.

The combination of large amount of data movement and large memory footprint, which we quantify in \secref{data_movement_and_roofline_model} and~\secref{memory_footprint} respectively, affects both software implementations running on commodity hardware, as well as custom hardware implementations. 
Commodity hardware (e.g., CPUs or GPUs) has a fixed amount of on-chip memory. The DP table might not fit into this on-chip memory, which introduces three inefficiencies: Data has to be moved a larger distance, which increases (1)~access latency and (2)~access energy~\citep{boroumand2018google}. (3)~The high amount of data movement puts high pressure on memory bandwidth, which is scarce when accessing data residing off-chip. This causes the entire application to become memory bandwidth-bound, thus wasting compute resources and achieving sub-optimal performance. Custom hardware implementations (e.g. ASICs) can have as large on-chip memory as needed, but such a large and high-bandwidth on-chip memory comes at the cost of a large chip area and power consumption (\citep{boroumand2021neuraledgedevices}).

Doing unnecessary work trivially wastes runtime and energy. In \secref{early_termination}, we identify the DP entries that are calculated needlessly by GenASM, and quantify how frequent they are.

\subsubsection{Roofline Model}
\label{sec:data_movement_and_roofline_model}
We use the \defn{roofline model}~\citep{seminal_roofline, applying_roofline} to visualize that GenASM has a large amount of data movement, and that its operational intensity (i.e., the number of operations per byte) is too low to saturate the compute resources of modern CPUs and GPUs. The roofline model plots the upper limit of achievable compute throughput for different operational intensities for a given processor. It consists of horizontal peak compute throughput rooflines, and sloped memory bandwidth rooflines.

\figref{roofline} shows the roofline plots for an Intel Xeon Gold 5118 CPU~\citep{xeon_gold_5118} and an NVIDIA A6000 GPU~\citep{A6000}, including their respective on-chip memory (\defn{shared memory} in CUDA), cache, and off-chip memory (\defn{global memory} in CUDA) bandwidths (drawn in shades of blue) and peak compute throughputs (draw in shades of green). GenASM's operational intensity is drawn in red. We derive the roofline parameters in \S8 of the supplementary materials.

\begin{figure}[h]
    \centering
    \includegraphics[width=\columnwidth,keepaspectratio]{figures/roofline_both.pdf}
    \caption{The roofline models of a) an Intel Xeon Gold 5118 CPU and b) an NVIDIA A6000 GPU.}
    \label{fig:roofline}
\end{figure}

From \figref{roofline} we make three observations. First, if the data resides off-chip, GenASM is heavily memory bandwidth-bound for a modern CPU and GPU. This is evidenced by the red (algorithm) and dark blue (off-chip memory bandwidth) lines intersecting far below the green (peak compute throughput) line.
Second, GenASM would no longer be memory bandwidth-bound if its computational intensity were $\geq$10\x higher, because then the red (algorithm) line would be shifted to the right and intersect with the dark blue (off-chip memory bandwidth) line above the green (peak compute throughput) line. The operational intensity could be increased by reducing GenASM's data movement.
Third, if the data resides in the fastest on-chip memory, GenASM \emph{can} reach peak compute throughput, even with the high amount of data movement in the baseline algorithm. This is evidenced by the red (algorithm) and light blue (L1/shared memory bandwidth) lines intersecting above the green (peak compute throughput) line. However, as we show in \secref{memory_footprint}, GenASM's memory footprint is too large for the typical capacity of such fast on-chip memories in commodity hardware, and building large enough on-chip memories is costly.

Based on these observations, we conclude that (1)~GenASM cannot saturate commodity hardware with computation, and (2)~data movement should be reduced to address this inefficiency.

\subsubsection{Memory Footprint} \label{sec:memory_footprint}
In this section, we demonstrate the overheads associated with GenASM's large memory footprint.%

We derive GenASM's working set memory footprint to be 96.5$KiB$ in \S9 of the supplementary materials. For comparison, the Intel Xeon Gold 5118 has 32$KiB$ of L1D cache per core~\citep{xeon_gold_5118} and NVIDIA's \defn{Ampere} GPU microarchitecture provides up to 99$KiB$ of high-bandwidth on-chip memory per GPU core (\defn{streaming multiprocessor}, \defn{SM} in CUDA)~\citep{cuda_guide}. Thus, one SM can hold the DP table for exactly one GenASM problem instance in its on-chip memory. One thread block of two warps (i.e., 2$\times$32 threads) can work on a single GenASM problem instance, but this does not saturate the compute resources in the SM. This is because modern GPUs are designed to alternate between executing \emph{multiple} independent instruction streams for the purpose of hiding the latency of instructions~\citep{lindholm2008nvidia}. Underutilization of the compute resources in an SM due to too few independent instruction streams is called \defn{low occupancy} and causes the unused computational resources to be wasted~\citep{cuda_guide}. Hence, the occupancy should be increased by working on multiple problem instances per SM. Multiple problem instances can fit into memory by \emph{either} reducing the memory footprint per problem instance, \emph{or} placing the DP tables into the GPU's off-chip memory, which has a much larger capacity. Our goal is the former, as we show in \secref{data_movement_and_roofline_model} that the latter is \emph{not} an efficient a solution due to the off-chip memory's limited bandwidth.

Custom hardware implementations (e.g., ASICs) can potentially have as large on-chip memory as needed. For example, the GenASM ASIC~\citep{cali2020genasm} uses scratchpads of 96.5~$KiB$ each to hold the DP tables. However, these scratchpads occupy 76\% of the total chip area and consume over 54\% of the chip power. This limits the performance achievable with a given chip area and power budget.

In summary, GenASM has a large memory footprint compared to typical on-chip memory capacities in commodity hardware, and while sufficiently large on-chip scratchpads can be designed for custom hardware implementations, it is costly to do so.

\subsection{Scrooge} \label{sec:optimizations}
We have shown in \secref{challenges} that GenASM has a high amount of data movement \emph{and} high memory footprint per problem instance. We have elaborated that this combination either limits performance (on commodity hardware), or requires expensive large on-chip memories (on custom hardware), both of which are undesirable.
Thus, our strategy is to reduce the GenASM algorithm's memory footprint as much as possible while introducing minimal computational overhead. We present three novel algorithmic improvements that collectively achieve a 24\x\ reduction in memory footprint, as well as a 12\x\ reduction in data movement from the memory that holds the DP table.

\subsubsection{Improvement~1 - Store Entries, not Edges (SENE)} \label{sec:SENE}
As we explain in \secref{traceback}, GenASM stores 3 bitvectors per entry of the table \var{R} to enable traceback. If we imagine a graph where the entries of \var{R} are nodes and the intermediate bitvectors are edges connecting their source and target entries, GenASM stores 3 ingoing edges for most nodes (\figref{SENE}).
We propose to trade off the majority of this memory footprint for a small increase in computation with the \defn{SENE} improvement. SENE regenerates the required edges on-demand during traceback from stored nodes (entries of table \var{R}) by applying the update rules in \algref{bitap} on requested neighbor entries.

\begin{figure}[h]
    \centering
    \includegraphics[width=\bitapfigwidth\columnwidth,keepaspectratio]{figures/sene.pdf}
    \caption{Per cell, GenASM stores three edges for traceback. \toolname with SENE stores only the DP entry itself instead; the needed edges are regenerated on the fly during traceback.}
    \label{fig:SENE}
\end{figure}

\textbf{Cost and benefits.}
Since traceback explores only a single path across the table \var{R}, only $O(\var{W})$ edges are regenerated, making the overhead of this extra computation small compared to computing the table of $O(\var{W}^2)$ entries. Storing \var{R} requires storing $65\times65$ entries of 64 bits each, for a total of $33,800B\approx 33kiB$. The previous memory footprint was $96.5kiB$ as derived in \secref{memory_footprint}, yielding a $\frac{96.5}{33} = 2.92 \approx 3\times$ improvement in memory footprint. Since each of these locations is still only written to once during the construction of \var{R}, SENE also reduces the data movement from the memory that holds the DP table by 3\x.

\subsubsection{Improvement~2 - Discard Entries not used by Traceback~(DENT)} \label{sec:DENT}
The windowing heuristic (\secref{windowing}) mandates that traceback covers only the first $\var{W}-\var{O}$ characters of each window. This means that traceback never reads the table entries of the last $\var{O}$ characters in each window. %

We propose to discard the entries that can never be reached by traceback, an improvement we call \defn{DENT}. These include the last $\var{O}$ columns of \var{R} and the last $\var{O-1}$ bits of every bitvector. The resulting DP table consists of $\var{W}-\var{O}+1$ columns, $\var{W}+1$ rows, and $\var{W}-\var{O}+1$ bits per entry. \figref{DENT} shows an example for $\var{W}$=4 and $\var{O}$=3, where \toolname stores only the leftmost 2 columns and leftmost 2 bits per entry, because traceback does not reach the rightmost 3 columns and rightmost 2 bits per entry.

\begin{figure}[h]
    \centering
    \includegraphics[width=\bitapfigwidth\columnwidth,keepaspectratio]{figures/dent.pdf}
    \caption{DENT exploits that the windowing heuristic stops traceback after the first $\var{W}-\var{O}$ edges are crossed (here $\var{W}$=4 and $\var{O}$=3). The area never reached by traceback can be discarded.}
    \label{fig:DENT}
\end{figure}

Due to the fixed word sizes and word alignment requirements of commodity hardware, the number of bits stored for each bitvector cannot be chosen freely. We show in \secref{throughput} that for a modern GPU $\var{O}=33$ achieves the best throughput results for $\var{W}=64$, because the stored bitvectors perfectly fit into a 32-bit word. In contrast, \citep{cali2020genasm} uses \var{O}=24 for its ASIC design, which we show to be suboptimal on commodity hardware. Note that increasing \var{O} \emph{improves} accuracy, see \secref{windowing} for an intuition and \secref{accuracy} for experimental results.

\textbf{Cost and benefits.}
DENT incurs two computational overheads: First, the bits to store have to be determined and extracted from the bitvectors. Second, increasing \var{O} from 24 to 33 means the algorithm makes 9 characters less progress per window.

By discarding the right half of each bitvector and the rightmost \var{O} columns of \var{R}, DENT improves the memory footprint by $\frac{\var{W}}{\var{W} - \var{O}+1}\times\frac{\var{W}+1}{\var{W} - \var{O}+1} = \frac{64}{32}\times\frac{65}{32} \approx 4\times$. We describe in \S1 of the supplementary materials how DENT can be extended to store only half the rows of \var{R} for a total 8\x memory footprint reduction.

By the same calculation, the number of writes to table \var{R} is reduced by approximately 4\x, assuming the forefront diagonal (marked red in \figref{bitap_example}) is kept in registers and communicated directly.

\subsubsection{Improvement~3 - Early Termination (ET)} \label{sec:early_termination}
The edit distance is determined by the highest row of \var{R} that contains a 0 in the most significant bit in the leftmost column. Traceback starts from this entry. Since entries are constructed from their north, north-east, and east neighbors, the traceback path can only go to the north, north-east, and east. It can \emph{never} go south. Thus, at no point do the rows of higher cost than $distance(\var{pattern}, \var{text})$ contain useful information for traceback (see \figref{early_termination} for an example). We propose building \var{R} row-wise, and terminating the algorithm early as soon as the most significant bit in the first entry of the current row is a 0.

\begin{figure}[h]
    \centering
    \includegraphics[trim={1.2cm 0 0 0},width=\bitapfigwidth\columnwidth,keepaspectratio]{figures/early_termination.pdf}
    \caption{The colored edges indicate the path taken by traceback for \var{W}=4 and \var{O}=0. Rows below the edit distance do not contain useful information for traceback. Thus, they do not need to be computed (Early~Termination).}
    \label{fig:early_termination}
\end{figure}

\textbf{Cost and benefits.}
Early Termination does not yield a constant factor improvement in either memory footprint or runtime: If $distance(\var{pattern}, \var{text}) = \var{W}$, we are not able to terminate early at all. However, typical input pairs incur fewer than \var{W} edits in a single window. For correct candidate pairs, the edit distance will be low, e.g. up to 15\% for long reads~\citep{alser2020technology}. Even uncorrelated random sequence pairs of length \var{W} over a 4-letter alphabet have an edit distance of at most $\frac{3}{4}$\var{W} on average, as we prove in \S10 of the supplementary materials. Thus, on average \toolname can skip at least 25\% of the entries of \var{R}, saving computation as well as data movement (see \secref{data_movement_and_roofline_model}).

\textbf{Conflict with intra-task parallelism.}
Recall from \secref{task_parallelism} that GenASM provides the option for intra-task parallelism.
Exploiting this parallelism requires the available processing elements to build \var{R} in a \emph{diagonal-wise} fashion, as shown in \figref{bitap_example}. However, as we describe in \secref{early_termination}, to make full use of Early~Termination, \var{R} should be built \emph{row-wise}. As a compromise, we implement Early~Termination in a diagonal-wise fashion in our GPU implementation. As in the row-wise version, construction on \var{R} stops as soon as the leftmost processing element finds a 0 in the most significant bit. Due to the diagonal-wise computation, the other processing elements have already computed several rows ahead at this point, i.e. done unnecessary work. For this reason, the benefit of Early~Termination is limited in intra-task parallel implementations, such as our GPU implementation, while being much more significant in row-wise implementations, such as our CPU implementation. We reaffirm these effects experimentally in \secref{sensitivity}.

\subsection{Implementation} \label{sec:implementation}
We implement C++ versions of our algorithm for x86 CPUs and NVIDIA GPUs. They are exposed as simple library functions for pairwise sequence alignment. Each improvement and implementation constant can be easily configured at compile time through preprocessor macros. The implementations, as well as baselines and evaluation scripts, are available at \gitrepo.

\textbf{CPU.}
The CPU version converts the input pairs to a two-bit-per-basepair encoding, but padded to 8 bits. Each thread works on a single pairwise alignment at a time and obtains sequence pairs from a global queue. During each call to the GenASM-DC subroutine, the thread calculates the DP table \var{R} in a row-wise fashion.

\textbf{GPU.}
The GPU version is implemented using CUDA~11.1~\citep{cuda_guide} and targets GPUs of compute capability~7.0 and higher~\citep{cuda_gpus}. The input sequence pairs are converted to a two-bit-per-basepair encoding and transferred to the GPU. Each thread block works on a single pairwise alignment at a time and obtains sequence pairs from a global queue. During each call to the GenASM-DC subroutine, the thread block calculates the DP table \var{R} in a diagonal-wise fashion, and each of the \var{W} threads in the thread block calculates a single column of \var{R}. Threads resolve their mutual data dependencies using warp shuffle instructions within a warp and using shared memory across warps. A single thread per warp executes the traceback operation. The size of the CIGAR string is not known ahead of time, hence it is stored as a linked list in global memory.


